\vspace{4pt}\noindent\textbf{Data generation.} A graph is grown from \rhval{const/founder_couples} founding couples. Each couple has between one and four children (probabilities \rhval{const/p_kids1}, \rhval{const/p_kids2}, \rhval{const/p_kids3} and \rhval{const/p_kids4} percent for one to four children), each child marries a newly created outsider with probability \rhval{const/p_marry} percent, and the new couple is queued, up to \rhval{const/max_gen} generations of couples, until exactly \rhval{const/people} people exist. Entity ids are randomly permuted so that ids carry no information. A graph has about \rhval{main/path-mp-2-hop/transductive/n_triples_g1/mean:0} triples (mean over seeds) spread almost evenly over the four relations (relation shares between \rhval{main/path-mp-2-hop/transductive/frac_uncle/mean:pct1}\% and \rhval{main/path-mp-2-hop/transductive/frac_parent/mean:pct1}\% on average).

\vspace{4pt}\noindent\textbf{Splits and the two settings.} Triples are split at the level of \emph{units} (\rhval{const/split_train}/\rhval{const/split_val}/\rhval{const/split_test} percent for train, validation and test); both directions of a sibling pair form one unit, so a held-out sibling triple has no mirror in the training data. Models are trained on the train+validation units of graph~1. \emph{Transductive} test: the held-out units of graph~1 (about \rhval{main/path-mp-2-hop/transductive/n_test_g1/mean:0} triples per seed), message graph = train+validation units of graph~1. \emph{Inductive} test: a second graph with \rhval{const/people} entities that were never seen, generated independently, split in the same proportions; message graph and fold-in data = its train+validation units; test = its held-out units (about \rhval{main/path-mp-2-hop/inductive/n_test_g2/mean:0} triples per seed). Both settings are therefore evaluated with the same fraction of the graph observed and the same number of test triples. A seed fixes both graphs, both splits, the initialisation and the minibatch order, so each seed is an independent graph pair and comparisons are paired by seed.

\vspace{4pt}\noindent\textbf{Systems and hyperparameters.} All learning systems are reimplemented in PyTorch (CPU, \rhval{const/threads} threads) in one script (\texttt{method/run.py}) with one evaluation function. TransE: dimension \rhval{const/embed_dim}, $L_1$ distance, margin $\gamma=\rhval{const/transe_margin}$, $k=\rhval{const/negs}$ negatives, Adam with learning rate \rhval{const/transe_lr:2}, batch size \rhval{const/transe_bs}, \rhval{const/transe_epochs} epochs; fold-in uses the same optimiser for \rhval{const/transe_epochs} epochs on the new graph. Path-MP: width $d=\rhval{const/pathmp_hidden}$, $L=2$, Adam with learning rate \rhval{const/pathmp_lr:3}, batches of \rhval{const/pathmp_bs} queries, \rhval{const/pathmp_epochs} epochs, gradient-norm clipping at \rhval{const/clip}.

\vspace{4pt}\noindent\textbf{Tuning.} Hyperparameters were chosen on the validation units of graph~1 for two separate graph seeds (\rhval{const/tune_seeds} seeds distinct from the evaluation seeds); the test units were never used. TransE received three rounds of search (margin in \{1,2,4,6,8,12,16\}, $L_1$ versus $L_2$ distance, dimension 64 versus 128, 200 versus 400 epochs), because the best margin was at the edge of the first grid; the optimum (margin 8, $L_1$, dimension 64, 200 epochs) lies inside the grid, and larger margins, the larger dimension (tried at margin 4 only) and more epochs (tried at margins 4, 8 and 12) did not improve validation MRR (all \rhval{count/tune_transe/runs} tuning runs are logged in group \texttt{tune\_transe}). Path-MP received a $2\times2$ grid (epochs 3 or 6, learning rate 0.003 or 0.01); 3 and 6 epochs were tied on validation and the cheaper 3 epochs were used. Path-MP therefore received a smaller tuning budget than TransE, which, if anything, favours the baseline (Section~\ref{sec:limitations}).

\vspace{4pt}\noindent\textbf{Statistics and runs.} \rhval{const/n_seeds} seeds (0--4). Tables report mean $\pm$ standard deviation over seeds; tests are paired $t$-tests over seeds from \texttt{rh compare} or from the analysis script (\texttt{method/analyze.py}, logged as a run). With five seeds these tests have little power for small effects and we do not interpret non-significant results as equivalence. Every run was executed with the \texttt{rh run} wrapper on a shared CPU machine; each run takes well below the per-run limit. The registry also contains failed rows that we did not hide: \rhval{const/n_failed_density} launches of the \texttt{sweep\_density} baselines (TransE+fold-in and the rule oracle) were corrupted by a variable clash in the launcher script and were re-run with the corrected script (the Path-MP rows of that sweep were unaffected), and one analysis-script launch failed on a registry control row before the script was fixed.
